% !TeX root = ../main.tex
% Rebuilt from: reliability-and-codes.tex, ch09-integrated-design.tex
\chapter{Reliability, Optimisation and Integrated Design}
\label{ch:integrated-design}

Each preceding chapter has produced quantities carrying uncertainty and has
passed them forward without examining it. This chapter addresses that
uncertainty directly: how it is described, how it propagates through a
calculation, how design codes accommodate it, and how a design is justified
when independent methods disagree and several objectives compete. The argument
has two parts. A design founded on a single set of parameter values has not
been tested against the range the ground actually presents, and the means of
testing it are neither elaborate nor costly. Integration, moreover, is a
process of decision rather than of averaging: where empirical, analytical and
numerical results diverge, the task is to establish which mechanism governs.

The argument has two halves. The first is that a design based on a single set
of parameter values has not been checked against the range the ground actually
presents, and that the tools for doing so are neither exotic nor expensive.
The second is that integration is a decision process rather than an averaging
process: where empirical, analytical and numerical methods diverge, the task
is to identify which mechanism governs, not to split the difference.

\section{Uncertainty and characteristic values}
\label{sec:uncertainty-and-characteristic-values}

\subsection{Design philosophies and uncertainty}
\label{sec:design-philosophies-and-uncertainty}

\paragraph{Three philosophies.}
Deterministic design adopts a single parameter set and demonstrates an
adequate factor of safety. It is simple, familiar and silent about how much of
the safety margin is absorbed by uncertainty rather than by genuine reserve.

The observational method adopts a most probable parameter set, establishes in
advance the range of behaviour that would be acceptable, prepares contingency
measures, and monitors to confirm which case is occurring. It is the natural
approach to underground works because the ground is progressively revealed,
and Eurocode~7 sets out the conditions it must satisfy: acceptable limits
defined beforehand, monitoring capable of detecting departure in time, and
contingency measures that can actually be implemented at that point.

Probabilistic design treats parameters as distributions and computes a
probability of unsatisfactory performance. It is the most informative and the
most demanding of data, and its value often lies less in the computed
probability than in what the calculation reveals about which parameters
matter.

The three are complementary rather than competing. Most underground projects
are designed deterministically, constructed observationally, and would benefit
from a probabilistic sensitivity study at the point where the deterministic
design is fixed.

\paragraph{Sources of uncertainty.}
Four sources are worth separating because they respond differently to effort.
\emph{Inherent spatial variability} is real variation in the ground and cannot
be reduced --- only characterised. \emph{Measurement error} arises in testing
and can be reduced by better technique. \emph{Transformation uncertainty}
arises when a measured quantity is converted to a design parameter through a
correlation, and it is frequently the largest contributor: the scatter in a
$\GSI$-to-modulus relation dwarfs the scatter in the underlying measurements.
\emph{Statistical estimation error} arises from estimating a distribution from
few samples, and it reduces with more testing.

The aleatory--epistemic distinction organises these. Aleatory uncertainty is
intrinsic randomness; epistemic uncertainty reflects lack of knowledge and is
reducible by investigation. Confusing them leads to investigation programmes
that cannot deliver what is asked of them, because no number of boreholes
removes the variability that is genuinely there.

\subsection{Parameter characterisation}
\label{sec:parameter-characterisation}

\paragraph{Statistical description.}
A parameter is described by its distribution, its mean and its coefficient of
variation. Typical coefficients of variation are well documented: intact rock
strength commonly 20 to 40 per cent, rock mass modulus higher, friction angle
considerably lower at 5 to 15 per cent, and unit weight lower still. These
figures are useful for sanity-checking a project data set, and a reported
coefficient of variation far below the published range usually means the data
have been filtered rather than that the ground is uniform.

Spatial correlation matters as much as scatter. Measurements close together
are more alike than measurements far apart, and the correlation length
describes the separation beyond which they become independent. It is the
reason that averaging over a failure surface reduces variance: a sliding
surface many correlation lengths long samples the full distribution, so the
mobilised strength approaches the mean, whereas a short surface may lie
entirely within one weak zone. Ignoring spatial averaging is conservative for
large failure surfaces and unconservative for local mechanisms.

\paragraph{Characteristic values.}
Eurocode~7 defines the characteristic value as a cautious estimate of the
value affecting the occurrence of the limit state. Every word of that
definition is doing work. \emph{Cautious} means on the safe side of the mean,
not the worst measurement. \emph{Affecting the occurrence of the limit state}
means the value relevant to the mechanism --- which may be a local minimum for
a small mechanism and a spatial average for a large one.

It follows that there is no universal percentile. A five per cent fractile is
a convention, not a definition, and applying it mechanically to a parameter
governing a large averaging mechanism is over-conservative, while applying it
to one governing a local mechanism may not be conservative at all.

\subsection{Limit state design and codes}
\label{sec:limit-state-design-and-codes}

\paragraph{Limit states in underground works.}
Ultimate limit states concern collapse: face instability, structural failure
of the lining, global instability. Serviceability limit states concern
function: convergence beyond the clearance envelope, settlement beyond a
third-party limit, cracking beyond watertightness requirements, leakage. In
urban tunnelling the serviceability states usually govern, and by a wide
margin.

Applying partial factors to underground works is awkward for a structural
reason. In a building, load and resistance are independent --- the wind does
not know the strength of the concrete. Underground, both derive from the same
ground: factoring the strength down also reduces the load the ground can shed,
and factoring the load up may change the mechanism. Eurocode~7's design
approaches handle this by applying factors at different points, and the choice
between them changes the answer. This is a genuine unresolved difficulty
rather than a matter of preference, and it is one reason tunnelling has
retained global factors of safety longer than other geotechnical disciplines.

\paragraph{Codes and guidance.}
No general design code covers tunnelling comprehensively. Geotechnical codes
provide the ground-related framework; structural and bridge codes cover
linings as concrete structures; and the specifically tunnelling guidance comes
from industry bodies --- the international tunnelling associations and the
national tunnelling societies --- which publish recommendations on segmental
lining design, sprayed concrete, risk management and fire.

The practical consequence is that a tunnel design is justified against a
composite of documents, none of which was written for the whole problem.
Stating which document governs which check is part of the design basis, and
projects that omit it discover the ambiguity during approval.

\section{Reliability and code frameworks}
\label{sec:reliability-and-code-frameworks}

\subsection{Reliability analysis}
\label{sec:reliability-analysis}

\paragraph{Reliability index and probability of failure.}
A performance function is defined as capacity minus demand, so that negative
values represent unsatisfactory performance. The reliability index $\beta$ is
the number of standard deviations between the mean of that function and zero,
and for a normally distributed performance function it maps directly to a
probability of failure: $\beta$ of about 2 corresponds to roughly two per
cent, 3 to about one in a thousand, and 4 to about three in a hundred
thousand.

Target values differ with consequence and duration. Temporary works during
construction, with people present but exposure short and monitoring active,
are designed to lower targets than a permanent lining that must perform for a
century with limited inspection. The computed probability should be read as a
notional and comparative quantity rather than as a prediction of frequency,
because it inherits every assumption in the input distributions.

\paragraph{Methods in practice.}
Four methods span the range of effort. First-order second-moment analysis
propagates means and variances through a linearised performance function and
is a spreadsheet calculation. The first-order reliability method locates the
most probable failure point and handles nonlinearity and non-normal
distributions properly. Point estimate methods evaluate the performance
function at a small number of chosen parameter combinations and require no
derivatives, which makes them well suited to problems whose evaluation is a
numerical model. Monte Carlo simulation samples the distributions directly and
is exact in the limit, at the cost of many evaluations.

Proportionality is the guide. A support design whose performance function is
an explicit expression justifies FORM or Monte Carlo. A design whose
evaluation is a finite element run justifies a point estimate method or a
surrogate model. Both are better than no uncertainty analysis, which is the
current default.

\subsection{Sensitivity and project risk}
\label{sec:sensitivity-and-project-risk}

\paragraph{Sensitivity analysis.}
Sensitivity analysis identifies which parameters govern the outcome, and it is
the highest-value uncertainty calculation available because it is cheap and
because it redirects effort. The characteristic result in rock engineering is
that the outcome is dominated by one or two parameters --- very often $\GSI$,
because rock mass strength depends on it exponentially and a ten-point change
can alter strength by 40 per cent, as the sensitivity study in
Chapter~\ref{ch:site-investigation} shows.

The consequence is a direct instruction to the investigation. If $\GSI$
governs and is poorly constrained, additional core logging and face mapping
buys more than additional strength testing. This is the same
value-of-information argument applied to investigation targeting in
Chapter~\ref{ch:ground-modelling}, and it is the mechanism by which
uncertainty analysis pays for itself.

\paragraph{From calculation to risk register.}
A computed probability becomes useful when it enters the project risk
framework, which speaks in likelihood and consequence rather than in $\beta$.
The translation requires the consequence to be assessed --- to life, to cost,
to programme, to third parties --- and the combination judged against
acceptance criteria, conventionally through an ALARP framing in which risk is
reduced until further reduction is grossly disproportionate to the benefit.

The insurance-driven codes of practice for tunnelling risk management make
this a contractual requirement rather than good practice, and the full
treatment is Chapter~\ref{ch:delivery-assets}.

\subsection{Worked example: uncertainty in a support design}
\label{sec:reliability-worked-example}

A representative calculation propagates uncertainty in $\GSI$, $\sci$ and
$\mi$ through the Hoek--Brown derivation of
Chapter~\ref{ch:rock-mass-parameters} to rock mass strength, then through a
support design to a factor of safety.

The instructive outcome is the comparison. A deterministic calculation at mean
parameters typically returns a comfortable factor of safety; the same
calculation with the distributions propagated returns a probability of
unsatisfactory performance that is small but not negligible, and a factor of
safety distribution whose lower tail is uncomfortably close to unity. The
deterministic answer is not wrong --- it is incomplete, because it says
nothing about how much of the margin is consumed by not knowing the ground.

The second outcome is the sensitivity ranking, which almost always places
$\GSI$ first by a wide margin. That result, obtained in an afternoon,
justifies redirecting investigation effort in a way no amount of additional
analysis at fixed parameters could.

\section{Integrating methods and objectives}
\label{sec:integrating-methods-and-objectives}

\subsection{The Design Workflow}

The four stages of underground geotechnical design form a logical sequence. Each stage builds on the outputs of the previous one:

\[
\text{Site Investigation} \xrightarrow{\text{params}} \text{Empirical} \xrightarrow{\text{support}} \text{Analytical} \xrightarrow{\text{CCM}} \text{Numerical}
\]

\paragraph{Stage 1: Site Investigation (Chapter~\ref{ch:site-investigation})}

\textbf{Shallow urban alignment:} published environmental statements and their geotechnical technical appendices, desk study of the surface constraints, and geological cross-sections along the alignment. Outputs: a two-layer ground model of a compressible surface unit over the host rock, parameter tables (Table~\ref{tab:srl_params}), groundwater conditions, and the third-party assets within the zone of influence.

\textbf{The deep pumped-storage scheme:} Geological mapping, borehole investigations, laboratory testing of intact sandstone/siltstone. Outputs: two rock mass conditions (Good: RMR\,=\,58, Q\,=\,8, GSI\,=\,60; Poor: RMR\,=\,34, Q\,=\,1.5, GSI\,=\,30), intact properties (UCS\,=\,72~MPa, $E_i$\,=\,26~GPa), joint set descriptions.

\textbf{Key decisions:}
\begin{itemize}[nosep]
 \item \textit{Shallow urban alignment:} is the surface unit thick enough along the alignment to influence tunnel response, and where does its contact with the host rock sit relative to the crown?
 \item \textit{The deep pumped-storage scheme:} Are the Good and Poor rock mass domains spatially separable along the alignment? Which domain governs support design?
 \item \textit{Both:} How reliable are the characteristic design values, given limited boreholes and core recovery issues?
\end{itemize}

\begin{table}[h!]
\centering
\small
\begin{tabular}{l l c c c c}
\toprule
\textbf{Unit} & \textbf{Description} & $\gamma$ (kN/m$^3$) & $c'$ (kPa) & $\phi'$ ($^\circ$) & $E$ (MPa) \\
\midrule
Upper & Compressible surface unit & 20 & 7 & 30 & 50 \\
Lower & Host rock unit & 21 & 20 & 28 & 150 \\
\bottomrule
\end{tabular}
\caption{Representative parameters for a shallow urban rail alignment.}
\label{tab:srl_params}
\end{table}

\begin{table}[h!]
\centering
\small
\begin{tabular}{l c c}
\toprule
\textbf{Parameter} & \textbf{Good rock (GSI = 60)} & \textbf{Poor rock (GSI = 30)} \\
\midrule
RMR & 58 & 34 \\
Q & 8 & 1.5 \\
UCS, $\sigma_{ci}$ (MPa) & 72 & 72 \\
Intact modulus, $E_i$ (GPa) & 26 & 26 \\
$m_i$ (assumed) & 17 & 17 \\
$m_b$ (HB) & 4.1 & 1.4 \\
$s$ (HB) & 0.012 & 0.0004 \\
$a$ (HB) & 0.503 & 0.522 \\
Rock mass strength, $\sigma_{cm}$ (MPa) & $\sim$12 & $\sim$2 \\
Rock mass modulus, $E_m$ (GPa) & $\sim$13 & $\sim$2.6 \\
MC equivalent $c'$ (MPa) & $\sim$3.0 & $\sim$0.6 \\
MC equivalent $\phi'$ ($^\circ$) & $\sim$45 & $\sim$30 \\
\bottomrule
\end{tabular}
\caption{Intact and derived rock mass parameters for the two design domains. Values are indicative and should be recomputed from the Hoek--Brown fit of Chapter~\ref{ch:ccm}.}
\label{tab:snowy_params}
\end{table}


\noindent\textbf{Link to later stages.}\enspace The shallow alignment parameters feed the numerical analysis of Chapter~\ref{ch:numerical-modelling}, and serve equally for empirical and analytical assessment of the station caverns. The deep scheme data feeds the empirical (Chapter~\ref{ch:ground-behaviour}) and analytical (Chapter~\ref{ch:ccm}) stages. Their reliability determines the confidence in all subsequent results---an error at Stage~1 propagates to every stage that follows.


\paragraph{Stage 2: Empirical Design (Chapter~\ref{ch:ground-behaviour})}

\textbf{Shallow urban alignment:} the host rock units fall in RMR and GSI ranges that permit empirical support design for both the running tunnels and the station caverns. RMR-based stand-up time can be assessed for the tunnel, and Q-system support can be designed for the underground station excavations.

\textbf{The deep pumped-storage scheme:} Failure mode identification using the decision tree and spalling equation of \citet{martin1999spalling}, and the Kaiser design chart. RMR-based stand-up time for the main access tunnel (span 9.5~m). Q-system support for the machine hall, transformer hall, and main access tunnel under both good and poor rock conditions.

\textbf{Key decisions:}
\begin{itemize}[nosep]
 \item \textit{Shallow urban alignment:} Which RMR class applies at the tunnel horizon? Is the stand-up time sufficient for the proposed advance rate?
 \item \textit{The deep pumped-storage scheme:} Is failure structurally controlled, stress-induced, or a combination? Does empirical support cover both Good and Poor domains?
 \item \textit{Both:} What preliminary support class should be carried forward into the analytical/numerical stages?
\end{itemize}


\noindent\textbf{Empirical methods assume structurally controlled failure.}\enspace RMR and Q-system were developed from case histories of jointed rock masses. They are not reliable for severe stress-induced spalling (which may apply to the deep pumped-storage scheme at 650~m depth). The decision tree helps identify when empirical methods are appropriate and when analytical/numerical methods are needed.


\paragraph{Stage 3: Analytical Design (Chapter~\ref{ch:ccm})}

\textbf{Shallow urban alignment:} Kirsch solutions~\citep{kirsch1898} can estimate elastic stress concentrations around the tunnel. For the shallow tunnel ($z_0 = 5$--25~m in the weathered host unit), the low stress means the response is likely elastic or mildly plastic. A simplified CCM analysis provides a cross-check against the numerical modelling results.

\textbf{The deep pumped-storage scheme:} Kirsch elastic stresses at various depths (100--600~m). Rock mass parameters ($m_b$, $s$, $a$, $\sigma_{cm}$, $E_m$, MC $c'$ and $\phi'$) for GSI\,=\,60 and GSI\,=\,30 at 650~m. Critical pressure, plastic zone radius, and the convergence--confinement equilibrium with the segmental lining.

\textbf{Key decisions:}
\begin{itemize}[nosep]
 \item \textit{Shallow urban alignment:} Is the analytical response elastic (no support driven by strength) or mildly plastic? Does CCM agree with numerical modelling?
 \item \textit{The deep pumped-storage scheme:} What is the required support capacity? Is the factor of safety adequate in the Poor rock domain?
 \item \textit{Both:} Where does the analytical solution disagree with empirical expectations, and why?
\end{itemize}

\begin{table}[h!]
\centering
\small
\begin{tabular}{l c c}
\toprule
\textbf{Parameter} & \textbf{GSI = 60} & \textbf{GSI = 30} \\
\midrule
Response & Plastic & Plastic \\
$R_p/a$ & 1.23 & 1.84 \\
$u_{\max}$ (mm) & 8.3 & 86.8 \\
$u_{s0}$ at $L=3$~m (mm) & 7.0 & 63.6 \\
\bottomrule
\end{tabular}
\caption{CCM results for the deep pumped-storage scheme main access tunnel at 650~m. The order-of-magnitude difference in displacement between GSI\,=\,60 and GSI\,=\,30 illustrates why rock mass characterisation dominates the design outcome---a 30-point change in GSI changes $u_{\max}$ by a factor of $\sim$10.}
\end{table}


\noindent\textbf{Analytical $\to$ Numerical.}\enspace The CCM provides displacement and FoS for a circular tunnel in hydrostatic stress. It cannot handle non-circular geometry, staged excavation, or layered ground. These limitations motivate numerical modelling.


\paragraph{Stage 4: Numerical Modelling (Chapter~\ref{ch:numerical-modelling})}

\textbf{Shallow urban alignment:} the empirical Gaussian trough~\citep{mair1996settlement} for surface settlement, with $V_L$ between 0.5 and 1.5\,per cent. Two-dimensional finite element models for the shallow ($z_0=5$~m) and deep ($z_0=25$~m) cases in the two-unit profile. Outputs: settlement troughs, displacement contours and lining forces, with the computed and empirical troughs overlaid.

\textbf{The deep pumped-storage scheme:} here the fourth stage is analytical verification rather than continuum modelling: the ground reaction, support characteristic and longitudinal deformation curves are computed and the equilibrium point located, by hand or with \companionname\ (Appendix~\ref{app:software}). That is adequate because the main access tunnel is close to circular and machine-driven, the stress field is nearly hydrostatic at depth, and the governing question is support capacity rather than ground--structure interaction. Two- or three-dimensional continuum analysis of the machine and transformer halls, which are non-circular and intersect one another, is the logical extension at detailed design.

\textbf{Key decisions:}
\begin{itemize}[nosep]
 \item \textit{Shallow urban alignment:} Does FEM confirm the Mair-predicted settlement trough? Are surface settlements within third-party damage criteria?
 \item \textit{The deep pumped-storage scheme:} Does CCM confirm that the TBM lining sits on a stable portion of the GRC? Is the FoS adequate in both domains?
 \item \textit{Both:} Where do numerical results contradict earlier stages, and what does that tell us about our assumptions?
\end{itemize}

\subsection{Complementarity and uncertainty}
\label{sec:complementarity-and-uncertainty}

\paragraph{How the Methods Complement Each Other}

\begin{table}[h!]
\centering
\small
\begin{tabularx}{\textwidth}{l X X X}
\toprule
\textbf{Purpose} & \textbf{Best method} & \textbf{Why} & \textbf{Limitation} \\
\midrule
Concept development & Empirical (RMR, Q) & Quick, requires minimal data, gives preliminary support & Assumes structurally controlled failure; no stress-strain detail \\
Identifying risks & Decision tree + SMR & Classifies failure mode; flags stress vs gravity vs squeezing & Qualitative; does not quantify FoS \\
Deformation assessment & Analytical (CCM) or Mair & Gives displacement estimate and FoS; fast parametric studies & Circular geometry, hydrostatic stress, simplified ground \\
Engineering decision & Numerical & Handles real geometry, layered ground, staged construction and lining forces & Requires validated parameters; mesh and boundary sensitivity; time-consuming \\
\bottomrule
\end{tabularx}
\caption{Which method for which purpose.}
\end{table}


\noindent\textbf{No single method is sufficient.}\enspace Empirical methods give a starting point. Analytical methods provide a cross-check and quantify ground-support interaction. Numerical methods handle the real complexity. The three approaches should be used together, not as alternatives. Disagreement between methods is informative---it highlights where assumptions break down.


\paragraph{Worked Example: Resolving Method Disagreement}

Suppose at the deep pumped-storage scheme main access tunnel in the Poor rock domain:
\begin{itemize}[nosep]
 \item \textbf{Empirical (Q-system):} Q\,=\,1.5, span 9.5~m $\to$ Category 5 support: systematic bolting + 50--90~mm fibre-reinforced shotcrete.
 \item \textbf{Analytical (CCM):} $u_{\max} = 86.8$~mm, $R_p/a = 1.84$, indicates a large plastic zone and significant squeezing potential. The Q-recommended support would likely be overstressed.
 \item \textbf{Numerical:} Support curve intersects GRC only after substantial convergence; lining safety factor marginal.
\end{itemize}

\noindent \textbf{Resolution.} The empirical recommendation is non-conservative for this stress/strength regime because the Q-system was calibrated against cases with lower squeezing. The engineer should: (i) adopt a stronger support class (e.g., yielding steel sets or thicker shotcrete with lattice girders), (ii) allow for convergence by using a flexible/yielding support philosophy, and (iii) flag the Poor rock domain as a high-risk zone requiring face mapping and probe drilling during construction. This is the judgement-driven synthesis that integrated design requires.

\paragraph{Uncertainty Through the Workflow}

Uncertainty enters at every stage and propagates forward:

\begin{enumerate}
 \item \textbf{Site investigation} (Chapter~\ref{ch:site-investigation}): RQD orientation bias, core loss, limited boreholes, spatial variability of GSI.
 \item \textbf{Empirical classification} (Chapter~\ref{ch:ground-behaviour}): Subjectivity in joint condition rating, SRF estimation, RMR/Q input errors ($\pm$10\% changes support category).
 \item \textbf{Analytical parameters} (Chapter~\ref{ch:ccm}): GSI $\pm$10 points changes $\sigma_{cm}$ by $\sim$40\%. MC fitting range affects $c'$ and $\phi'$.
 \item \textbf{Numerical modelling} (Chapter~\ref{ch:numerical-modelling}): Mesh sensitivity, boundary effects, constitutive model choice, $K_0$ assumption, volume loss assumption in Mair method.
\end{enumerate}


\noindent\textbf{Project-specific uncertainties: a shallow urban tunnel.}\enspace \begin{itemize}[nosep]
 \item The depth of the contact between the two units is inferred from adjacent boreholes, not measured at the tunnel location.
 \item The host rock parameters ($c'=20$~kPa, $\phi'=28^\circ$) are equivalent continuum values; the rock mass itself is more variable than a single pair implies.
 \item The shallow case ($z_0=5$~m, $D=7.2$~m) is at the limit of the Mair method's validity ($z_0 < D$).
 \item Groundwater conditions were simplified (dry or piezometric) --- real conditions may involve transient pore pressures.
 \item $K_0$ was assumed, not measured; this significantly affects horizontal stresses and lining forces.
\end{itemize}


\noindent\textbf{Uncertainty as part of the recommendation.}\enspace An integrated design is incomplete unless it identifies the uncertainties that govern its conclusion, and these normally arise at several stages of the workflow rather than at one. For each, the design record should state how the recommendation would change if the assumption proved wrong, what additional data would reduce the uncertainty, and at what cost. Stated in that form, uncertainty ceases to be a qualification attached to the conclusion and becomes part of it.


\subsection{Carbon as a design objective}
\label{sec:integrated-carbon}

\paragraph{Extending the objective set.}
\label{sec:objective-set}
Factor of safety, displacement, settlement, cost and constructability are the
conventional objectives, and a design is normally justified by demonstrating
that each is satisfied. Adding carbon (Chapter~\ref{ch:sustainability}) does
not introduce a foreign criterion: the design already determines carbon
through geometry, support timing and excavation method, and the only question
is whether that influence is exercised deliberately or by default.

Making it explicit changes the calculation in one respect. The conventional
objectives are mostly constraints --- satisfied or not --- whereas carbon and
cost are quantities to be minimised. A design problem with several constraints
and two competing objectives has a different structure from one with
constraints alone, and it does not have a unique answer.

\paragraph{Pareto trade-offs.}
\label{sec:integrated-pareto}
Where objectives compete, the set of interest is the Pareto set: those designs
for which no objective can be improved without worsening another. Designs
outside it are simply inferior and can be discarded; designs within it
represent genuine trades that require a judgement rather than a calculation.

Reporting the set rather than a single optimum is the honest outcome, because
it exposes what the client is actually being asked to accept --- so much
carbon for so much reserve, so much cost for so much convergence. A single
recommended design conceals the trade and, in practice, conceals the weighting
the designer applied without saying so.

For problems where each evaluation is a numerical model, generating the set
directly is impractical, and surrogate models fitted to a modest number of
runs are the usual route. Genetic and other population-based algorithms are
appropriate once a surrogate exists; simple parametric comparison over a
handful of candidate designs is often enough and should be tried first.

\paragraph{Controlled deformation as a design lever.}
\label{sec:controlled-deformation-lever}
The clearest illustration that carbon and geotechnical performance are coupled
rather than independent is the support timing decision. Permitting more ground
convergence before support is mobilised moves the equilibrium point down the
ground reaction curve, reducing the support pressure the lining must carry and
therefore its thickness and reinforcement --- a material and carbon reduction
obtained through mechanics rather than through substitution.

The limits are real and must be stated with the saving. Convergence consumes
the clearance envelope, so the excavated profile must accommodate it.
Excessive convergence loosens the rock mass and permanently reduces its
strength, which moves the ground reaction curve itself. In shallow ground,
convergence appears at the surface as settlement and runs into the third-party
limits of Chapter~\ref{ch:settlement-damage}. And a design that relies on a
particular installation distance has made an assumption about construction
that must appear in the specification.

Within those limits the lever is genuine, and it is the example worth carrying
from this chapter: the carbon reduction available from understanding the
mechanics exceeds what is available from specifying a different concrete.

\section{Workflow structure and project linkage}
\label{sec:workflow-structure-and-project-linkage}

\paragraph{Example Workflow Structure}

An integrated design normally carries the structure below, adapted to the analysis actually performed. The stages are a sequence of dependence, not a timetable: each consumes what the previous produced, and an error early in the chain propagates through everything after it.

\begin{table}[h!]
\centering
\small
\begin{tabularx}{\textwidth}{c l X X}
\toprule
\textbf{Stage} & \textbf{Activity} & \textbf{Shallow urban alignment} & \textbf{Deep pumped-storage scheme} \\
\midrule
1 & Site investigation & Published statements, desk study, two-unit parameters & Geological mapping, RMR/Q/GSI, intact properties \\
2 & Empirical design & RMR and Q for station and running tunnel & Failure mode, RMR and Q for caverns and access tunnel \\
3 & Analytical design & Kirsch, simplified convergence--confinement cross-check & Kirsch at depth, Hoek--Brown and Mohr--Coulomb parameters, convergence--confinement \\
4 & Numerical & Empirical trough and finite element model & Analytical verification, continuum analysis as extension \\
5 & Integration & \multicolumn{2}{c}{Workflow synthesis and review of the assumption chain} \\
\bottomrule
\end{tabularx}
\caption{Two contrasting projects carried through the same design sequence. The stage that governs differs; the sequence does not.}
\end{table}

\paragraph{Linking the Two Projects}

Both projects follow the same four-stage design workflow, but they emphasise different aspects---and \textit{the emphasis follows from the engineering problem}, not from arbitrary choice:

\textbf{A shallow urban alignment.} A shallow urban tunnel in weathered rock/soil. The governing engineering concern is surface settlement and its impact on existing buildings, services, and roads. Because the critical design question is \textit{how much will the ground above the tunnel move}, the decision-making tools are the Mair method and FEM ---i.e., Stage 1 (site investigation) and Stage 4 (numerical modelling) dominate. Empirical classification of the AF rock units and analytical stress/displacement estimates still apply, but they inform rather than govern the design.

\textbf{The deep pumped-storage scheme.} A deep tunnel and cavern system in hard rock at high stress. The governing engineering concern is rock mass failure mode and support capacity. Because the critical design question is \textit{will the rock mass fail around the opening and what support will hold it}, the decision-making tools are the failure-mode decision tree, RMR/Q support, and the CCM---i.e., Stage 2 (empirical) and Stage 3 (analytical) dominate. The project still begins with site investigation data, and the convergence--confinement analysis serves as the verification stage.


\noindent\textbf{Both projects demonstrate the complete workflow.}\enspace The key integration insight is that the same four-stage sequence applies to \textit{any} underground project, regardless of depth, rock type, or excavation method. Shallow tunnels in weak ground and deep caverns in hard rock both require site investigation $\to$ empirical classification $\to$ analytical assessment $\to$ numerical verification. The difference lies in which stage is most critical and which tools are most appropriate at each stage---and that difference is driven by the engineering question, not by preference.


\section{Updating the design with observation}
\label{sec:bayesian-updating}

The preceding sections treat uncertainty as a fixed property of the design:
distributions are assigned, propagated, and a reliability index results. That
is the right picture before construction and the wrong one after it begins.
From the first round of excavation the project generates evidence about the
same parameters it assumed, and a design method that cannot absorb that
evidence discards the most relevant data it will ever have.

This section supplies the machinery. It is what the observational method of
Chapter~\ref{ch:movement-monitoring} performs informally, what the ground
model of Chapter~\ref{ch:ground-modelling} requires in order to be kept
current, and what the value-of-information argument
(Section~\ref{sec:value-of-information}) presupposes.

\subsection{The updating relation}
\label{sec:updating-relation}

Let $\theta$ be a parameter of interest --- a rock mass modulus, a volume
loss, a friction angle --- described before construction by a probability
density $p(\theta)$ expressing what the investigation established. Let $z$ be
an observation whose value depends on $\theta$ through a model with a known
error, so that $p(z \mid \theta)$ states how probable that observation would
be for each candidate value. The revised description of $\theta$ after seeing
$z$ is
\[
p(\theta \mid z) \;=\; \frac{p(z \mid \theta)\, p(\theta)}{p(z)} ,
\]
in which the denominator is a normalising constant and carries no information
about $\theta$. The three terms have names worth keeping straight: the prior
$p(\theta)$, the likelihood $p(z \mid \theta)$ --- read as a function of
$\theta$ for the $z$ actually observed, not as a distribution over $z$ --- and
the posterior $p(\theta \mid z)$.

For the common case in which prior and measurement error are both normal, the
result is closed-form and instructive. With a prior of mean $\mu_0$ and
standard deviation $\sigma_0$, and a single observation $z$ of standard
deviation $\sigma_z$, the posterior is normal with
\[
\mu_1 = \frac{\sigma_z^2 \mu_0 + \sigma_0^2 z}{\sigma_0^2 + \sigma_z^2},
\qquad
\frac{1}{\sigma_1^2} = \frac{1}{\sigma_0^2} + \frac{1}{\sigma_z^2} .
\]
Three properties follow directly and are worth more than the algebra. The
posterior mean is a weighted average of prediction and measurement, each
weighted by the other's variance, so a precise measurement dominates a vague
prior and a vague measurement barely moves a confident prior. Precisions add,
so the posterior is always at least as sharp as either input: an observation
cannot increase uncertainty about a parameter, though it can move the mean a
long way. And the update is sequential --- yesterday's posterior is today's
prior --- so a drive can be updated round by round without reprocessing its
history.

The practical difficulty is rarely the arithmetic. It is that $\sigma_z$ must
include the error of the model connecting $\theta$ to $z$, not merely the
precision of the instrument. A convergence array reading to a tenth of a
millimetre does not deliver a modulus to comparable precision, because the
relation between the two carries the assumptions of the analysis that
produced it. Treating instrument precision as measurement precision is the
commonest error in applied updating, and it produces a posterior far sharper
than the evidence supports.

\subsection{Updating in construction}
\label{sec:updating-in-construction}

What is worth updating is the small set of parameters that both govern the
design and are observable during construction. Deformation modulus qualifies:
it controls convergence and lining load, and convergence is measured on every
project. Volume loss qualifies in soft ground, being both the governing
settlement input and directly observable from the trough. In-situ stress
qualifies where spalling depth is recorded
(Chapter~\ref{ch:in-situ-stress}). Intact strength usually does not, being
better determined in the laboratory than inferred from tunnel behaviour.

Two failures recur. The first is updating a parameter that the observation
does not identify. Convergence depends on stress, modulus and strength
together, and a single convergence record cannot separate them: any of several
combinations reproduces it, so an update attributing all of the discrepancy to
modulus is arbitrary. The remedy is to update the identifiable combination, or
to bring in a second independent observation --- plastic zone extent from
extensometers, or spalling depth --- so that the joint likelihood
discriminates.

The second is the prior that will not move. A prior stated with unrealistically
small variance, usually because the design team is defending its
characterisation rather than describing it, produces posteriors that ignore the
evidence. The test is honest and simple: state before construction what
observation would change the design, and how far the parameter would have to
move. A prior for which no realistic observation changes anything is not a
prior but a decision already taken.

Used properly the loop is unremarkable in operation. A prediction is made and
recorded with its uncertainty; the excavation is monitored; the posterior is
computed; the support class ahead of the face is set from the posterior rather
than the original assumption. The discipline lies entirely in recording the
prediction before the observation, which is what distinguishes updating from
hindsight.

\subsection{The value of an observation}
\label{sec:value-of-observation}

The same machinery answers a question posed but not resolved in
Section~\ref{sec:value-of-information}: what is a measurement worth before it
is taken.

An observation has value only through the decision it might change. If the
support class is the same for every outcome the measurement could produce,
the measurement is worth nothing however large the uncertainty it resolves.
The formal statement follows from the updating relation: for each possible
outcome $z$, compute the posterior, select the decision that is best under
that posterior, and evaluate its expected cost. Averaging those costs over the
prior predictive distribution of $z$ gives the expected cost with the
measurement; the difference from the expected cost of the best decision under
the prior alone is the expected value of the information. A measurement is
worth taking when that difference exceeds its cost.

Two consequences are worth stating because they are frequently resisted. The
value depends on the decision, so the same borehole is worth different amounts
to different designs, and an investigation cannot be optimised without knowing
what will be done with it. And information about a parameter the design is
insensitive to has no value whatever its geological interest --- which is the
formal version of the argument made qualitatively in
Chapter~\ref{ch:ground-modelling}, and the reason investigation targeting
belongs to the designer rather than to the driller.

The calculation is rarely performed in full, and it need not be. Its value in
practice is that it forces the three questions that determine the answer:
what decision is at stake, what outcomes could the measurement produce, and
would any of them change what is done. A programme that cannot answer those
three is not yet an investigation programme
(\citealt{baecher2003reliability} sets out the formal treatment).
